\documentclass[12pt]{article}
\usepackage[UTF8]{inputenc}
\usepackage{xeCJK}
\setCJKmainfont{Sarabun}
\usepackage{enumitem}
\usepackage{geometry}
\geometry{a4paper, margin=2.5cm}
\usepackage{setspace}
\onehalfspacing

\begin{document}

\begin{center}
{\Large\textbf{ฟิสิกส์อะตอม}}\\[6pt]
{\normalsize วิชาฟิสิกส์ ม.ปลาย บทที่ 25}\\[4pt]
{\footnotesize ทั้งหมด 25 ข้อ}
\end{center}
\hrule
\bigskip

\section*{คำถาม in class}
\begin{enumerate}[label=\arabic*., itemsep=0.5\baselineskip, topsep=0.5\baselineskip, leftmargin=*]
\item แบบจำลองอะตอมของดอลตัน (Dalton) บอกว่าอะตอมเป็นอย่างไร?
\begin{enumerate}[label=\alph*)., itemsep=0.3\baselineskip, topsep=0.3\baselineskip]
\item ก\ 
\begin{minipage}[t]{0.9\linewidth}
\raggedright มีประจุบวกและประจุลบผสมกันทั่ว
\end{minipage}
\item ข\ 
\begin{minipage}[t]{0.9\linewidth}
\raggedright เป็นทรงกลมแข็งที่ไม่แบ่งแยกได้
\end{minipage}
\item ค\ 
\begin{minipage}[t]{0.9\linewidth}
\raggedright มีนิวเคลียร์ที่ประกอบด้วยโปรตอนและนิวตรอน
\end{minipage}
\item ง\ 
\begin{minipage}[t]{0.9\linewidth}
\raggedright มีอิเล็กตรอนวิ่งรอบนิวเคลียร์เหมือนระบบสุริยะ
\end{minipage}
\end{enumerate}
\vspace{-0.3\baselineskip}
\begin{small}
\textit{คำตอบ: ข}
\end{small}

\item แบบจำลองอะตอมของทอมสัน (Thomson) เรียกอีกอย่างว่าอะไร?
\begin{enumerate}[label=\alph*)., itemsep=0.3\baselineskip, topsep=0.3\baselineskip]
\item ก\ 
\begin{minipage}[t]{0.9\linewidth}
\raggedright แบบจำลองนิวเคลียร์
\end{minipage}
\item ข\ 
\begin{minipage}[t]{0.9\linewidth}
\raggedright แบบจำลองลูกกลมโป้ยปุ้น (Plum pudding model)
\end{minipage}
\item ค\ 
\begin{minipage}[t]{0.9\linewidth}
\raggedright แบบจำลองระบบสุริยะ
\end{minipage}
\item ง\ 
\begin{minipage}[t]{0.9\linewidth}
\raggedright แบบจำลองควอนตัม
\end{minipage}
\end{enumerate}
\vspace{-0.3\baselineskip}
\begin{small}
\textit{คำตอบ: ข}
\end{small}

\item การทดลองของรัทเทอร์ฟอร์ด (Rutherford) ใช้อนุภาคอะไรยิงไปที่แผ่นทองคำ?
\begin{enumerate}[label=\alph*)., itemsep=0.3\baselineskip, topsep=0.3\baselineskip]
\item ก\ 
\begin{minipage}[t]{0.9\linewidth}
\raggedright อิเล็กตรอน
\end{minipage}
\item ข\ 
\begin{minipage}[t]{0.9\linewidth}
\raggedright โปรตอน
\end{minipage}
\item ค\ 
\begin{minipage}[t]{0.9\linewidth}
\raggedright อนุภาคแอลฟา (α)
\end{minipage}
\item ง\ 
\begin{minipage}[t]{0.9\linewidth}
\raggedright นิวตรอน
\end{minipage}
\end{enumerate}
\vspace{-0.3\baselineskip}
\begin{small}
\textit{คำตอบ: ค}
\end{small}

\item จากการทดลองของรัทเทอร์ฟอร์ด ข้อใดไม่ถูกต้องเกี่ยวกับผลการทดลอง?
\begin{enumerate}[label=\alph*)., itemsep=0.3\baselineskip, topsep=0.3\baselineskip]
\item ก\ 
\begin{minipage}[t]{0.9\linewidth}
\raggedright อนุภาคแอลฟาส่วนใหญ่ทะลุแผ่นทองคำตรง
\end{minipage}
\item ข\ 
\begin{minipage}[t]{0.9\linewidth}
\raggedright มีอนุภาคแอลฟาบางตัวกระเจิงในมุมใหญ่
\end{minipage}
\item ค\ 
\begin{minipage}[t]{0.9\linewidth}
\raggedright อนุภาคแอลฟาทุกตัวสะท้อนกลับทิศทางเดิม
\end{minipage}
\item ง\ 
\begin{minipage}[t]{0.9\linewidth}
\raggedright มีนิวเคลียร์ขนาดเล็กรวมประจุบวกอยู่ตรงกลาง
\end{minipage}
\end{enumerate}
\vspace{-0.3\baselineskip}
\begin{small}
\textit{คำตอบ: ค}
\end{small}

\item แบบจำลองอะตอมของโบร์ (Bohr model) ใช้กับอะตอใดได้ดีที่สุด?
\begin{enumerate}[label=\alph*)., itemsep=0.3\baselineskip, topsep=0.3\baselineskip]
\item ก\ 
\begin{minipage}[t]{0.9\linewidth}
\raggedright อะตอมไฮโดรเจน
\end{minipage}
\item ข\ 
\begin{minipage}[t]{0.9\linewidth}
\raggedright อะตอมฮีเลียม
\end{minipage}
\item ค\ 
\begin{minipage}[t]{0.9\linewidth}
\raggedright อะตอมคาร์บอน
\end{minipage}
\item ง\ 
\begin{minipage}[t]{0.9\linewidth}
\raggedright อะตอมออกซิเจน
\end{minipage}
\end{enumerate}
\vspace{-0.3\baselineskip}
\begin{small}
\textit{คำตอบ: ก}
\end{small}

\item ระดับพลังงานของอิเล็กตรอนในอะตอมมีลักษณะอย่างไร?
\begin{enumerate}[label=\alph*)., itemsep=0.3\baselineskip, topsep=0.3\baselineskip]
\item ก\ 
\begin{minipage}[t]{0.9\linewidth}
\raggedright มีค่าเท่ากันทุกระดับ
\end{minipage}
\item ข\ 
\begin{minipage}[t]{0.9\linewidth}
\raggedright มีค่าเป็นเส้นตรงเพิ่มขึ้นต่อเนื่อง
\end{minipage}
\item ค\ 
\begin{minipage}[t]{0.9\linewidth}
\raggedright มีค่าเฉพาะที่ไม่ต่อเนื่อง (เป็นควอนตัม)
\end{minipage}
\item ง\ 
\begin{minipage}[t]{0.9\linewidth}
\raggedright เป็นค่าต่ำสุดเสมอ
\end{minipage}
\end{enumerate}
\vspace{-0.3\baselineskip}
\begin{small}
\textit{คำตอบ: ค}
\end{small}

\item เมื่ออิเล็กตรอนในอะตอมไฮโดรเจนกระโดดจากระดับ n=3 ลงมายังระดับ n=1 จะเกิดอะไรขึ้น?
\begin{enumerate}[label=\alph*)., itemsep=0.3\baselineskip, topsep=0.3\baselineskip]
\item ก\ 
\begin{minipage}[t]{0.9\linewidth}
\raggedright ดูดกลืนโฟตอน
\end{minipage}
\item ข\ 
\begin{minipage}[t]{0.9\linewidth}
\raggedright ปล่อยโฟตอน
\end{minipage}
\item ค\ 
\begin{minipage}[t]{0.9\linewidth}
\raggedright อิเล็กตรอนหลุดจากอะตอม
\end{minipage}
\item ง\ 
\begin{minipage}[t]{0.9\linewidth}
\raggedright ไม่มีการเปลี่ยนแปลงพลังงาน
\end{minipage}
\end{enumerate}
\vspace{-0.3\baselineskip}
\begin{small}
\textit{คำตอบ: ข}
\end{small}

\item สเปกตรัมเส้น (Line spectrum) ของอะตอมเกิดจากอะไร?
\begin{enumerate}[label=\alph*)., itemsep=0.3\baselineskip, topsep=0.3\baselineskip]
\item ก\ 
\begin{minipage}[t]{0.9\linewidth}
\raggedright การเปล่งแสงของวัตถุร้อนทั้งหมด
\end{minipage}
\item ข\ 
\begin{minipage}[t]{0.9\linewidth}
\raggedright การเปลี่ยนระดับพลังงานของอิเล็กตรอนในอะตอม
\end{minipage}
\item ค\ 
\begin{minipage}[t]{0.9\linewidth}
\raggedright การกระจายของแสงขาว
\end{minipage}
\item ง\ 
\begin{minipage}[t]{0.9\linewidth}
\raggedright การสะท้อนของแสงบนผิววัตถุ
\end{minipage}
\end{enumerate}
\vspace{-0.3\baselineskip}
\begin{small}
\textit{คำตอบ: ข}
\end{small}

\item แบบจำลองอะตอมของโบร์ กับ ของรัทเทอร์ฟอร์ด แตกต่างกันอย่างไร?
\begin{enumerate}[label=\alph*)., itemsep=0.3\baselineskip, topsep=0.3\baselineskip]
\item ก\ 
\begin{minipage}[t]{0.9\linewidth}
\raggedright โบร์บอกว่ามีนิวเคลียร์ รัทเทอร์ฟอร์ดบอกว่าไม่มี
\end{minipage}
\item ข\ 
\begin{minipage}[t]{0.9\linewidth}
\raggedright โบร์บอกว่าอิเล็กตรอนโคจรโดยไม่แผ่รังสี รัทเทอร์ฟอร์ดบอกว่าอิเล็กตรอนเคลื่อนที่แบบใดก็ได้
\end{minipage}
\item ค\ 
\begin{minipage}[t]{0.9\linewidth}
\raggedright โบร์ใช้ได้กับทุกอะตอม รัทเทอร์ฟอร์ดใช้ได้กับไฮโดรเจน
\end{minipage}
\item ง\ 
\begin{minipage}[t]{0.9\linewidth}
\raggedright รัทเทอร์ฟอร์ดบอกว่ามีนิวเคลียร์ โบร์เพิ่มเงื่อนไขว่าอิเล็กตรอนอยู่ในระดับพลังงานคงที่
\end{minipage}
\end{enumerate}
\vspace{-0.3\baselineskip}
\begin{small}
\textit{คำตอบ: ง}
\end{small}

\item สเปกตรัมเส้นของไฮโดรเจนมีชุดสเปกตรัมที่สำคัญกี่ชุด?
\begin{enumerate}[label=\alph*)., itemsep=0.3\baselineskip, topsep=0.3\baselineskip]
\item ก\ 
\begin{minipage}[t]{0.9\linewidth}
\raggedright 1 ชุด
\end{minipage}
\item ข\ 
\begin{minipage}[t]{0.9\linewidth}
\raggedright 2 ชุด
\end{minipage}
\item ค\ 
\begin{minipage}[t]{0.9\linewidth}
\raggedright 3 ชุด
\end{minipage}
\item ง\ 
\begin{minipage}[t]{0.9\linewidth}
\raggedright 4 ชุด
\end{minipage}
\end{enumerate}
\vspace{-0.3\baselineskip}
\begin{small}
\textit{คำตอบ: ง}
\end{small}

\item ระดับพลังงานของอิเล็กตรอนในอะตอมไฮโดรเจนหาได้จากสูตรใด?
\begin{enumerate}[label=\alph*)., itemsep=0.3\baselineskip, topsep=0.3\baselineskip]
\item ก\ 
\begin{minipage}[t]{0.9\linewidth}
\raggedright E\_n = -hcR/n
\end{minipage}
\item ข\ 
\begin{minipage}[t]{0.9\linewidth}
\raggedright E\_n = -hcR/n²
\end{minipage}
\item ค\ 
\begin{minipage}[t]{0.9\linewidth}
\raggedright E\_n = hcR/n²
\end{minipage}
\item ง\ 
\begin{minipage}[t]{0.9\linewidth}
\raggedright E\_n = n²hcR
\end{minipage}
\end{enumerate}
\vspace{-0.3\baselineskip}
\begin{small}
\textit{คำตอบ: ข}
\end{small}

\item ชุด Balmer ในสเปกตรัมของไฮโดรเจน เป็นเส้นที่อิเล็กตรอนกระโดดลงมายังระดับใด?
\begin{enumerate}[label=\alph*)., itemsep=0.3\baselineskip, topsep=0.3\baselineskip]
\item ก\ 
\begin{minipage}[t]{0.9\linewidth}
\raggedright n = 1
\end{minipage}
\item ข\ 
\begin{minipage}[t]{0.9\linewidth}
\raggedright n = 2
\end{minipage}
\item ค\ 
\begin{minipage}[t]{0.9\linewidth}
\raggedright n = 3
\end{minipage}
\item ง\ 
\begin{minipage}[t]{0.9\linewidth}
\raggedright n = 4
\end{minipage}
\end{enumerate}
\vspace{-0.3\baselineskip}
\begin{small}
\textit{คำตอบ: ข}
\end{small}

\item เมื่ออิเล็กตรอนดูดกลืนโฟตอนที่มีพลังงานเท่ากับผลต่างระดับพลังงานสองระดับ จะเกิดอะไรขึ้น?
\begin{enumerate}[label=\alph*)., itemsep=0.3\baselineskip, topsep=0.3\baselineskip]
\item ก\ 
\begin{minipage}[t]{0.9\linewidth}
\raggedright อิเล็กตรอนจะหลุดจากอะตอม
\end{minipage}
\item ข\ 
\begin{minipage}[t]{0.9\linewidth}
\raggedright อิเล็กตรอนจะกระโดดขึ้นไประดับพลังงานที่สูงกว่า
\end{minipage}
\item ค\ 
\begin{minipage}[t]{0.9\linewidth}
\raggedright อิเล็กตรอนจะอยู่ที่เดิม
\end{minipage}
\item ง\ 
\begin{minipage}[t]{0.9\linewidth}
\raggedright นิวเคลียร์จะเปลี่ยนแปลง
\end{minipage}
\end{enumerate}
\vspace{-0.3\baselineskip}
\begin{small}
\textit{คำตอบ: ข}
\end{small}

\item หลอดฟลูออเรสเซนต์ทำงานอย่างไร?
\begin{enumerate}[label=\alph*)., itemsep=0.3\baselineskip, topsep=0.3\baselineskip]
\item ก\ 
\begin{minipage}[t]{0.9\linewidth}
\raggedright ใช้ไฟฟ้าจ่ายพลังงานให้อิเล็กตรอนในอะตอมปรอท ซึ่งปล่อยรังสี UV ไปกระตุ้นผงเรืองแสงให้เปล่งแสงขาว
\end{minipage}
\item ข\ 
\begin{minipage}[t]{0.9\linewidth}
\raggedright ใช้ความร้อนทำให้วัตถุเรืองแสง
\end{minipage}
\item ค\ 
\begin{minipage}[t]{0.9\linewidth}
\raggedright ใช้ปฏิกิริยาเคมีโดยตรง
\end{minipage}
\item ง\ 
\begin{minipage}[t]{0.9\linewidth}
\raggedright ใช้แม่เหล็กกระตุ้นอิเล็กตรอน
\end{minipage}
\end{enumerate}
\vspace{-0.3\baselineskip}
\begin{small}
\textit{คำตอบ: ก}
\end{small}

\item ทำไมแบบจำลองอะตอมของรัทเทอร์ฟอร์ดจึงไม่สามารถอธิบายเสถียรภาพของอะตอมได้?
\begin{enumerate}[label=\alph*)., itemsep=0.3\baselineskip, topsep=0.3\baselineskip]
\item ก\ 
\begin{minipage}[t]{0.9\linewidth}
\raggedright เพราะไม่มีนิวเคลียร์
\end{minipage}
\item ข\ 
\begin{minipage}[t]{0.9\linewidth}
\raggedright เพราะอิเล็กตรอนเคลื่อนที่เป็นวงโคจรจะแผ่รังสีแม่เหล็กไฟฟ้าทำให้สูญเสียพลังงานและตกลงนิวเคลียร์
\end{minipage}
\item ค\ 
\begin{minipage}[t]{0.9\linewidth}
\raggedright เพราะนิวเคลียร์มีขนาดใหญ่เกินไป
\end{minipage}
\item ง\ 
\begin{minipage}[t]{0.9\linewidth}
\raggedright เพราะไม่มีอิเล็กตรอน
\end{minipage}
\end{enumerate}
\vspace{-0.3\baselineskip}
\begin{small}
\textit{คำตอบ: ข}
\end{small}

\item พลังงานไอออไนเซชัน (ionization energy) ของอะตอมไฮโดรเจนคือเท่าใด?
\begin{enumerate}[label=\alph*)., itemsep=0.3\baselineskip, topsep=0.3\baselineskip]
\item ก\ 
\begin{minipage}[t]{0.9\linewidth}
\raggedright 3.4 eV
\end{minipage}
\item ข\ 
\begin{minipage}[t]{0.9\linewidth}
\raggedright 13.6 eV
\end{minipage}
\item ค\ 
\begin{minipage}[t]{0.9\linewidth}
\raggedright 1.51 eV
\end{minipage}
\item ง\ 
\begin{minipage}[t]{0.9\linewidth}
\raggedright 0.85 eV
\end{minipage}
\end{enumerate}
\vspace{-0.3\baselineskip}
\begin{small}
\textit{คำตอบ: ข}
\end{small}

\item ความยาวคลื่นของเส้น H\_α (H-alpha) ในสเปกตรัมไฮโดรเจน ซึ่งเป็นเส้นแรกของชุด Balmer เกิดจากการกระโดดจากระดับใด?
\begin{enumerate}[label=\alph*)., itemsep=0.3\baselineskip, topsep=0.3\baselineskip]
\item ก\ 
\begin{minipage}[t]{0.9\linewidth}
\raggedright n=2 → n=1
\end{minipage}
\item ข\ 
\begin{minipage}[t]{0.9\linewidth}
\raggedright n=3 → n=2
\end{minipage}
\item ค\ 
\begin{minipage}[t]{0.9\linewidth}
\raggedright n=4 → n=2
\end{minipage}
\item ง\ 
\begin{minipage}[t]{0.9\linewidth}
\raggedright n=5 → n=2
\end{minipage}
\end{enumerate}
\vspace{-0.3\baselineskip}
\begin{small}
\textit{คำตอบ: ข}
\end{small}

\item สเปกตรัมดูดกลืน (absorption spectrum) ของไฮโดรเจนจะมีลักษณะอย่างไรเมื่อให้แสงขาวผ่านไอไฮโดรเจน?
\begin{enumerate}[label=\alph*)., itemsep=0.3\baselineskip, topsep=0.3\baselineskip]
\item ก\ 
\begin{minipage}[t]{0.9\linewidth}
\raggedright แถบมืดในช่วงแสงที่ตามองเห็นทั้งหมด
\end{minipage}
\item ข\ 
\begin{minipage}[t]{0.9\linewidth}
\raggedright เส้นมืดเฉพาะบริเวณที่ตรงกับการกระโดดของอิเล็กตรอน
\end{minipage}
\item ค\ 
\begin{minipage}[t]{0.9\linewidth}
\raggedright ไม่มีแถบมืดเลย
\end{minipage}
\item ง\ 
\begin{minipage}[t]{0.9\linewidth}
\raggedright แถบมืดเฉพาะในช่วง UV
\end{minipage}
\end{enumerate}
\vspace{-0.3\baselineskip}
\begin{small}
\textit{คำตอบ: ข}
\end{small}

\item ปรากฏการณ์ฟอสฟอเรสเซนซ์ (phosphorescence) ต่างจากฟลูออเรสเซนซ์ (fluorescence) อย่างไร?
\begin{enumerate}[label=\alph*)., itemsep=0.3\baselineskip, topsep=0.3\baselineskip]
\item ก\ 
\begin{minipage}[t]{0.9\linewidth}
\raggedright ฟอสฟอเรสเซนซ์เกิดทันที ฟลูออเรสเซนซ์เกิดช้า
\end{minipage}
\item ข\ 
\begin{minipage}[t]{0.9\linewidth}
\raggedright ฟอสฟอเรสเซนซ์เกิดช้า (มีการเก็บพลังงานไว้ก่อนแล้วค่อยปล่อย)
\end{minipage}
\item ค\ 
\begin{minipage}[t]{0.9\linewidth}
\raggedright ฟอสฟอเรสเซนซ์ใช้รังสี UV ฟลูออเรสเซนซ์ใช้แสงขาว
\end{minipage}
\item ง\ 
\begin{minipage}[t]{0.9\linewidth}
\raggedright ไม่ต่างกัน
\end{minipage}
\end{enumerate}
\vspace{-0.3\baselineskip}
\begin{small}
\textit{คำตอบ: ข}
\end{small}

\item ข้อใดไม่ใช่ข้อจำกัดของแบบจำลองอะตอมของโบร์?
\begin{enumerate}[label=\alph*)., itemsep=0.3\baselineskip, topsep=0.3\baselineskip]
\item ก\ 
\begin{minipage}[t]{0.9\linewidth}
\raggedright ใช้ได้เฉพาะกับอะตอมที่มีอิเล็กตรอน 1 ตัว
\end{minipage}
\item ข\ 
\begin{minipage}[t]{0.9\linewidth}
\raggedright อธิบายสเปกตรัมของอะตอมฮีเลียมไม่ได้
\end{minipage}
\item ค\ 
\begin{minipage}[t]{0.9\linewidth}
\raggedright อธิบายปรากฏการณ์โฟโตอิเล็กทริกได้อย่างสมบูรณ์
\end{minipage}
\item ง\ 
\begin{minipage}[t]{0.9\linewidth}
\raggedright อธิบายโครงสร้างละเอียดของสเปกตรัมเส้นไม่ได้
\end{minipage}
\end{enumerate}
\vspace{-0.3\baselineskip}
\begin{small}
\textit{คำตอบ: ค}
\end{small}

\item อิเล็กตรอนในอะตอมไฮโดรเจนอยู่ที่ระดับ n=2 ต้องการพลังงานเท่าใดจึงจะหลุดจากอะตอม?
\begin{enumerate}[label=\alph*)., itemsep=0.3\baselineskip, topsep=0.3\baselineskip]
\item ก\ 
\begin{minipage}[t]{0.9\linewidth}
\raggedright 13.6 eV
\end{minipage}
\item ข\ 
\begin{minipage}[t]{0.9\linewidth}
\raggedright 10.2 eV
\end{minipage}
\item ค\ 
\begin{minipage}[t]{0.9\linewidth}
\raggedright 3.4 eV
\end{minipage}
\item ง\ 
\begin{minipage}[t]{0.9\linewidth}
\raggedright 1.89 eV
\end{minipage}
\end{enumerate}
\vspace{-0.3\baselineskip}
\begin{small}
\textit{คำตอบ: ข}
\end{small}

\item ความยาวคลื่นของเส้นแรกในชุด Lyman (Ly-α) มีค่าเท่าใด? (กำหนด R = 1.097×10⁷ m⁻¹)
\begin{enumerate}[label=\alph*)., itemsep=0.3\baselineskip, topsep=0.3\baselineskip]
\item ก\ 
\begin{minipage}[t]{0.9\linewidth}
\raggedright 91.2 nm
\end{minipage}
\item ข\ 
\begin{minipage}[t]{0.9\linewidth}
\raggedright 121.6 nm
\end{minipage}
\item ค\ 
\begin{minipage}[t]{0.9\linewidth}
\raggedright 656.3 nm
\end{minipage}
\item ง\ 
\begin{minipage}[t]{0.9\linewidth}
\raggedright 486.1 nm
\end{minipage}
\end{enumerate}
\vspace{-0.3\baselineskip}
\begin{small}
\textit{คำตอบ: ข}
\end{small}

\item อิเล็กตรอนในอะตอมไฮโดรเจนอยู่ที่สถานะพื้น (ground state) ดูดกลืนโฟตอนที่มีพลังงาน 12.75 eV อิเล็กตรอนจะกระโดดไปถึงระดับใด?
\begin{enumerate}[label=\alph*)., itemsep=0.3\baselineskip, topsep=0.3\baselineskip]
\item ก\ 
\begin{minipage}[t]{0.9\linewidth}
\raggedright n = 2
\end{minipage}
\item ข\ 
\begin{minipage}[t]{0.9\linewidth}
\raggedright n = 3
\end{minipage}
\item ค\ 
\begin{minipage}[t]{0.9\linewidth}
\raggedright n = 4
\end{minipage}
\item ง\ 
\begin{minipage}[t]{0.9\linewidth}
\raggedright n = 5
\end{minipage}
\end{enumerate}
\vspace{-0.3\baselineskip}
\begin{small}
\textit{คำตอบ: ค}
\end{small}

\item หลังจากอิเล็กตรอนถูกกระตุ้นขึ้นสู่ระดับพลังงานสูงแล้ว มันจะกลับลงมาสู่สถานะพื้นได้กี่แบบ?
\begin{enumerate}[label=\alph*)., itemsep=0.3\baselineskip, topsep=0.3\baselineskip]
\item ก\ 
\begin{minipage}[t]{0.9\linewidth}
\raggedright 1 แบบเท่านั้น
\end{minipage}
\item ข\ 
\begin{minipage}[t]{0.9\linewidth}
\raggedright 2 แบบเท่านั้น
\end{minipage}
\item ค\ 
\begin{minipage}[t]{0.9\linewidth}
\raggedright หลายแบบ (หลายเส้นทาง)
\end{minipage}
\item ง\ 
\begin{minipage}[t]{0.9\linewidth}
\raggedright ไม่สามารถกลับลงมาได้
\end{minipage}
\end{enumerate}
\vspace{-0.3\baselineskip}
\begin{small}
\textit{คำตอบ: ค}
\end{small}

\item ข้อใดเป็นเหตุผลว่าทำไมสเปกตรัมเส้นของแต่ละธาตุจึงเฉพาะเจาะจงต่อธาตุนั้น?
\begin{enumerate}[label=\alph*)., itemsep=0.3\baselineskip, topsep=0.3\baselineskip]
\item ก\ 
\begin{minipage}[t]{0.9\linewidth}
\raggedright เพราะแต่ละธาตุมีมวลต่างกัน
\end{minipage}
\item ข\ 
\begin{minipage}[t]{0.9\linewidth}
\raggedright เพราะระดับพลังงานของอิเล็กตรอนเฉพาะเจาะจงต่อธาตุนั้น
\end{minipage}
\item ค\ 
\begin{minipage}[t]{0.9\linewidth}
\raggedright เพราะนิวเคลียร์ของแต่ละธาตุมีขนาดเท่ากัน
\end{minipage}
\item ง\ 
\begin{minipage}[t]{0.9\linewidth}
\raggedright เพราะอิเล็กตรอนทุกตัวมีความเร็วเท่ากัน
\end{minipage}
\end{enumerate}
\vspace{-0.3\baselineskip}
\begin{small}
\textit{คำตอบ: ข}
\end{small}

\end{enumerate}
\end{document}